\documentclass[12pt]{amsart}
\usepackage{amsmath,amssymb,amsthm,hyperref,xcolor,tabularx}

% ── Placeholder convention ─────────────────────────────
% \PLACEHOLDER{text} marks items to fill after decisions
\newcommand{\PLACEHOLDER}[1]{{\color{red}\textbf{[#1]}}}
% ── Status markers ─────────────────────────────────────
\newcommand{\DONE}[1]{{\color{blue}\textbf{[DONE: #1]}}}
\newcommand{\UPDATED}[1]{{\color{teal}\textbf{[UPDATED 2026-04-21: #1]}}}

\newtheorem{theorem}{Theorem}[section]
\newtheorem{conjecture}[theorem]{Conjecture}
\newtheorem{corollary}[theorem]{Corollary}
\newtheorem{lemma}[theorem]{Lemma}
\newtheorem{definition}[theorem]{Definition}
\newtheorem{remark}[theorem]{Remark}
\newtheorem{example}[theorem]{Example}

\title[PCF Arithmetic Stratification]{%
A Complete Arithmetic Stratification of\\
Polynomial Continued Fractions}

\author{papanokechi}
\address{Yokohama, Japan}
\urladdr{https://github.com/papanokechi}
\email{\PLACEHOLDER{email}}

\keywords{polynomial continued fractions,
arithmetic stratification, spectral classification,
obstruction hierarchy, PSL$_2(\mathbb{Z})$,
experimental mathematics, Wallis recurrence,
rational contamination, PSLQ}

\subjclass[2020]{11A55, 11J70, 11Y65, 11F06, 11B37}

\thanks{%
  Supplementary data, SIARC relay logs, and AEAL
  evidence records:
  \href{https://github.com/papanokechi/pcf-research}{%
  github.com/papanokechi/pcf-research}.
  Individual stratum results are developed in the
  companion papers
  \cite{Rat,Log,Alg,Obstruction,Spectral,Desert,Khinchin};
  this paper synthesizes them into a unified framework.
  Contamination methodology:
  \href{https://github.com/papanokechi/pcf-rational-contamination}{%
  github.com/papanokechi/pcf-rational-contamination}.}

\begin{document}
\maketitle

% ══════════════════════════════════════════════════════
% STATUS DASHBOARD — update this block as decisions arrive
% ══════════════════════════════════════════════════════
%
% Submission pipeline (as of 2026-04-20):
%
%  [1] JNTH-D-26-00480     JNT          — under review
%  [2] ExpMath 260636324   Exp Math     — under review
%  [3] ExpMath 266535114   Exp Math     — under review
%  [4] ExpMath 268704746   Exp Math     — under review
%  [5] ExpMath 266999523   Exp Math     — under review
%  [6] ExpMath 267262713   Exp Math     — under review
%  [7] ExpMath 264514392   Exp Math     — under review
%  [8] ExpMath 261945835   Exp Math     — SUBMITTED 2026-04-20
%       (pcf_rational_contamination_2026.pdf)
%  [17] NON-110708    Nonlinearity  — SUBMITTED 2026-04-21
%       (vquad_resurgence_SUBMISSION.pdf)
%
% TRIGGER: begin drafting body text after
% ≥2 ExpMath decisions received.
% JNT decision on [1] = spine signal.
% Nonlinearity decision on [17] = F2 signal.
%
% ══════════════════════════════════════════════════════

\begin{abstract}
We develop a complete arithmetic stratification of
the space $\mathcal{F}(d,D)$ of degree-$d$ polynomial
continued fractions with integer coefficients bounded
by $D$. The space decomposes into five mutually
exclusive strata --- Rational (\textbf{Rat}),
Logarithmic (\textbf{Log}), Algebraic (\textbf{Alg}),
Transcendental (\textbf{Trans}), and Desert
(\textbf{Des}) --- each characterized by an explicit
algebraic obstruction within a four-tier hierarchy
connecting PCF arithmetic to the action of
$\mathrm{PSL}_2(\mathbb{Z})$ and the Takeuchi
classification of arithmetic Fuchsian groups.

\medskip
The \textbf{Rat} stratum is fully characterized:
a PCF converges to a rational limit if and only if
$a(1)=0$ (trivial zero: $K=b(0)$) or $a(k)=0$ for
some $k\geq 2$ (finite termination), accounting for
$43/1000$ families ($4.3\%$) in a random survey
across degrees $d=2,\ldots,6$ \cite{Rat}.
The \textbf{Log} stratum is characterized by a
Frobenius obstruction of Catalan--logarithmic type
\cite{Log}. The \textbf{Desert} stratum accounts for
$95.7\%$ of random families at precision
\texttt{dps}$=150$; its non-emptiness is witnessed
by precision-controlled PSLQ null results \cite{Desert,Khinchin}.
The obstruction hierarchy spine is established via
the spectral 5-tuple
$\mathrm{Spec}(K)=(d,\Lambda,\Delta,\rho,\tau)$
and the four-tier framework of \cite{Obstruction}.

\medskip
A completeness conjecture --- that every convergent
PCF in $\mathcal{F}(d,D)$ falls into exactly one
stratum --- is supported by a $1000$-family
experimental survey with zero unclassified families.

\smallskip
\noindent\textit{AI assistance disclosure.}
Results were discovered through the SIARC
(Self-Iterating Analytic Relay Chain) multi-model
relay pipeline (iterations 1--15+).
Manuscript preparation involved Claude
(\texttt{claude-sonnet-4-6}, Anthropic) and
GitHub Copilot (Microsoft). All mathematical
results were conceived, directed, and validated
by the author.
\end{abstract}

% ══════════════════════════════════════════════════════
\section{Introduction}
\label{sec:intro}
% ══════════════════════════════════════════════════════

A \emph{polynomial continued fraction} (PCF) of
degree $d$ is an expression
\[
  K \;=\; b(0) + \mathbf{K}_{n=1}^{\infty}
  \frac{a(n)}{b(n)},
  \qquad a(n), b(n) \in \mathbb{Z}[n],
  \quad \deg a, \deg b \leq d.
\]
The space $\mathcal{F}(d,D)$ consists of all such
PCFs with $\|a\|_\infty, \|b\|_\infty \leq D$.

\subsection{The stratification problem}

\begin{quote}
\textit{Given a convergent PCF
$K \in \mathcal{F}(d,D)$, what determines the
arithmetic nature of its limit, and what obstructs
identification when the limit is transcendental?}
\end{quote}

\subsection{Main results}

\begin{enumerate}
\item \textbf{(Stratification,
  \S\S\ref{sec:rat}--\ref{sec:desert})}
  $\mathcal{F}(d,D)$ decomposes as
  $\mathrm{Rat} \sqcup \mathrm{Log} \sqcup
  \mathrm{Alg} \sqcup \mathrm{Trans} \sqcup
  \mathrm{Des}$,
  with explicit membership criteria.

\item \textbf{(Rat stratum, \S\ref{sec:rat})}
  $K \in \mathrm{Rat}$ iff $a(1)=0$ (giving $K=b(0)$)
  or $\exists\,k\geq 2: a(k)=0$.
  \DONE{proved, submitted as \cite{Rat},
  ExpMath 261945835, 2026-04-20}

\item \textbf{(Rational contamination,
  \S\ref{sec:rat})}
  The enriched PSLQ basis
  $\{1,K,\ln 2, K\ln 2,\ln^2 2\}$ produces
  false-positive transcendental identities at
  rate $\sim\!4.3\%$ from Rat-stratum families.
  Pre-screening protocol eliminates both
  contamination classes.
  \DONE{proved and documented in \cite{Rat}}

\item \textbf{(Log stratum, \S\ref{sec:log})}
  \PLACEHOLDER{sharpen after ExpMath decision
  on 266999523}

\item \textbf{(Obstruction hierarchy,
  \S\ref{sec:obstruction})}
  Four-tier filtration: Wallis $\to$ Frobenius
  $\to$ Spectral $\to$ Fuchsian.
  \PLACEHOLDER{upgrade Conjecture to Theorem
  after JNT decision on JNTH-D-26-00480}

\item \textbf{(Desert dominance,
  \S\ref{sec:desert})}
  Under uniform sampling with $D=4$,
  $957/1000$ families are Desert at
  \texttt{dps}$=150$.
  \UPDATED{empirical result from
  1000-family classification pass,
  2026-04-20}

\item \textbf{(Genuine ln(2) null result,
  \S\ref{sec:desert})}
  A corrected search over 800 families
  ($d=2,\ldots,5$, $D=4$) using stripped
  basis $\{1,\ln 2,\ln^2 2\}$ with
  rationality pre-screen finds zero genuine
  ln(2) identities at residual $<10^{-40}$.
  \UPDATED{null result confirmed 2026-04-20}

\item \textbf{(Completeness conjecture,
  \S\ref{sec:completeness})}
  Every convergent $K\in\mathcal{F}(d,D)$
  falls into exactly one stratum.
  Evidence: 1000-family survey,
  0 unclassified.
  \UPDATED{Both-stratum artifact resolved:
  all 41 multi-basis hits confirmed rational,
  2026-04-20}
\end{enumerate}

% ══════════════════════════════════════════════════════
\section{The Spectral 5-Tuple}
\label{sec:spectral}
% ══════════════════════════════════════════════════════

\begin{definition}[Spectral class]
For $K \in \mathcal{F}(d,D)$ convergent, define
\[
  \mathrm{Spec}(K) \;=\;
  (d,\,\Lambda,\,\Delta,\,\rho,\,\tau),
\]
where $d$ is the polynomial degree; $\Lambda$ is
the asymptotic Poincar\'e exponent of $(Q_n)$;
$\Delta$ is the discriminant of the connection
matrix of the associated rank-2 recurrence;
$\rho$ is the PSLQ spectral density; and
$\tau \in \{0,1,2,3\}$ is the obstruction
tier following \cite{Obstruction}
(\S\ref{sec:obstruction}); the pre-tier
$\tau = -1$ denotes Rat-stratum families
where the Wallis recurrence collapses before
any hypergeometric structure is reached.
\end{definition}

\PLACEHOLDER{Insert Poincar\'e characteristic
lemma with $\Lambda$ formula from JNT paper
after JNTH-D-26-00480 decision.}

% ══════════════════════════════════════════════════════
\section{The Rational Stratum}
\label{sec:rat}
% ══════════════════════════════════════════════════════

\begin{theorem}[Rat stratum characterization
  {\cite{Rat}}]
\label{thm:rat}
Let $K \in \mathcal{F}(d,D)$ be convergent.
\begin{enumerate}
\item[(i)] If $a(1)=0$, then $K=b(0)$
  (trivial zero mechanism).
\item[(ii)] If $a(k)=0$ for some $k\geq 2$
  and $a(n)\neq 0$ for $1\leq n<k$,
  then $K=P_{k-1}/Q_{k-1}\in\mathbb{Q}$
  (finite termination mechanism).
\end{enumerate}
Together, (i) and (ii) account for all
$43/1000$ rational-limit families in the
survey, with zero residual cases.
\end{theorem}

\begin{remark}[Rational contamination in
  PSLQ searches]
The enriched basis $\{1,K,\ln 2,K\ln 2,\ln^2 2\}$
expresses rational $K$ via spurious factored
relations $(c_1 K+c_0)(1+\ln 2)=0$, producing
false-positive transcendental identity claims
at rate $\sim\!4{-}9\%$ depending on degree.
Pre-screening via $[1,K]$ PSLQ before any
enriched-basis call eliminates this artifact.
Full methodology in \cite{Rat}.
\end{remark}

\begin{corollary}[Rational density]
Under uniform sampling from $\{-D,\ldots,D\}$,
\[
  P(K \in \mathrm{Rat}) \;\approx\;
  \frac{1}{\sqrt{2\pi(d+1)D(D+1)/3}}
  \;+\; O(D^{-(d+2)/2}),
\]
empirically $3$--$8.5\%$ at $d=2,\ldots,6$,
$D=4$.
\end{corollary}

% ══════════════════════════════════════════════════════
\section{The Logarithmic Stratum}
\label{sec:log}
% ══════════════════════════════════════════════════════

\PLACEHOLDER{Full section to be written after
ExpMath decisions on 266999523 and 266535114.}

\medskip\noindent\textbf{Structural framework
(from \cite{Obstruction}, available now):}

The Log stratum decomposes into two
geometrically distinct sub-types corresponding
to Tiers~0 and~1 of the obstruction hierarchy:

\begin{itemize}
\item \textbf{Tier-0 Log} (modular,
  generic transcendental).
  Seven verified families
  (\textnormal{log-}$k=2$,
  \textnormal{pi-}$m=0,\ldots,5$)
  with limits lying on the modular orbifold
  $X(1) = \mathrm{PSL}_2(\mathbb{Z})\backslash
  \mathfrak{H}^*$.
  The limits are generically transcendental
  but not CM.
  Spectral coordinates: $\Lambda=3$,
  $\tau=0$.

\item \textbf{Tier-1 Log} (arithmetic cuspidal,
  logarithmic).
  Prototype: Catalan, with Frobenius residual
  $3/(2n)$ and fitted signature $(\infty,1,4)$.
  The limit involves logarithmic terms
  incompatible with a clean Gauss identity.
  Spectral coordinates: $\Lambda$ modified by
  logarithmic Frobenius growth, $\tau=1$.
\end{itemize}

\begin{remark}[Log stratum inaccessibility under
  uniform polynomial sampling]
A directed search over 1500 polynomial families
($d=2,3,4$, $D=4$, $N=500$ per degree) with
pre-filter threshold $10^{-6}$ against target
constants $\{\ln 2, \ln 3, \ln 5, G, \pi/4\}$
found zero pre-filter passes: no random polynomial
PCF limit came within $10^{-6}$ of any target.
The classical Log examples (Catalan, $\ln 2$ via
Gauss CF) use non-polynomial partial numerators
outside $\mathcal{F}(d,D)$. This confirms that
$\mu(\mathrm{Log}) = 0$ under the uniform measure
on $\mathcal{F}(d,D)$ and that Log families require
structural construction beyond uniform sampling.
\UPDATED{confirmed 2026-04-20, iteration 16}
\end{remark}

\noindent\PLACEHOLDER{Insert: Frobenius
obstruction lemma (Catalan-logarithmic type),
482 irrational constants theorem,
$4/\pi$ Casoratian identity, and Wallis--PCF
arithmetic sector bifurcation. Cite \cite{Log}
and \cite{Alg} from accepted versions.}

% ══════════════════════════════════════════════════════
\section{The Algebraic Stratum}
\label{sec:alg}
% ══════════════════════════════════════════════════════

\PLACEHOLDER{Write after ExpMath decisions on
266535114 and 268704746.}

% ══════════════════════════════════════════════════════
\section{The Four-Tier Obstruction Hierarchy}
\label{sec:obstruction}
% ══════════════════════════════════════════════════════

\begin{definition}[Schwarz triangle data]
For a $d=1$ PCF admitting a Gauss hypergeometric
witness $\,_2F_1(a,b;c;z)$, the Schwarz triangle
parameters are
\[
  \lambda = |c-1|,\quad
  \mu = |a-b|,\quad
  \nu = |a+b-c|,
\]
giving triangle order triple $(\infty,\mu,\nu)$
when $\lambda=0$.
\end{definition}

\begin{theorem}[\textnormal{PSL}$_2(\mathbb{Z})$
  theorem {\cite{Obstruction}, Theorem~1}]
\label{thm:psl2z}
Every arithmetic $d=1$ polynomial continued
fraction family in the current verified taxonomy
has hypergeometric Schwarz triangle group
commensurable with $\mathrm{PSL}_2(\mathbb{Z})$
and uniformizes the classical modular curve
$X(1)$.

The arithmetic families are exactly those with
$c=1$ (forcing $\lambda=0$) and
$(\mu,\nu)=(2,3)$: seven verified families
(\textnormal{log-}$k=2$,
\textnormal{pi-}$m=0,\ldots,5$),
all governed by the same
$\,_2F_1(-1-i,-1+i;1;z)$ at different
evaluation points $z_0$ on the same orbifold.
The condition $c=1$ is necessary but not
sufficient; the sharper pair $(\mu,\nu)=(2,3)$
distinguishes the arithmetic sector.
\end{theorem}

\begin{definition}[Obstruction tiers
  {\cite{Obstruction}, Section~5}]
\label{def:tiers}
\begin{itemize}
\item \textbf{Tier~0 (Modular).}
  Arithmetic; signature $(\infty,2,3)$;
  $c=1$, $(\mu,\nu)=(2,3)$.
  Proved obstruction: the PCF evaluation point
  $z_0$ is generically transcendental and
  non-CM; CM occurs only at isolated modularly
  distinguished arguments.
  Examples: \textnormal{log-}$k=2$,
  \textnormal{pi-}$m=0,\ldots,5$.

\item \textbf{Tier~1 (Arithmetic cuspidal).}
  Arithmetic in origin but Schwarz triangle
  degenerates cuspidally; fitted signature
  $(\infty,1,4)$; Frobenius-log residual
  $3/(2n)$.
  Two independent proofs: analytic (Frobenius
  residual, \cite{Spectral}) and geometric
  (cuspidal triangle degeneration,
  \cite{Obstruction}).
  Example: Catalan.

\item \textbf{Tier~2 (Non-arithmetic Fuchsian).}
  Gauss $\,_2F_1$ witness exists but triangle
  order triple $(\infty,q,r)$ absent from
  Takeuchi's list \cite{Takeuchi}.
  Proved obstruction: no Shimura-curve CM
  interpretation exists at the group-theoretic
  level (Lemma~1 of \cite{Obstruction}).
  Examples: \textnormal{log-}$k=3,4,5,6,7,8,10$;
  some have irrational orders
  $\sqrt{3}$, $\sqrt{5}$.

\item \textbf{Tier~3 (Irregular Heun).}
  No Gauss $\,_2F_1$ witness; governing
  differential equation has irregular
  singularity at $\infty$; outside the
  Fuchsian triangle-group world.
  Proved obstruction (\cite{Obstruction},
  Proposition~1): neither Takeuchi's list
  nor the Tier-0 modular criterion applies.
  Example: $V_{\mathrm{quad}}\approx
  1.19737\ldots$, confirmed as confluent-Heun
  type at CM discriminant $-11$.
\end{itemize}
\end{definition}

\begin{remark}[Poincar\'e exponent and tiers]
For a $d=1$ PCF in Tier~0 or~1, the
Poincar\'e exponent of the Wallis denominator
sequence satisfies $\Lambda = \mathrm{Re}(\nu)$
where $\nu=|a+b-c|$ is the Schwarz parameter.
For Tier-0 families $\nu=3$, giving $\Lambda=3$.
For Tiers~2 and~3, $\Lambda$ is irrational or
undefined within the Gauss framework.
\end{remark}

\begin{remark}[$V_{\mathrm{quad}}$ as Tier-3
  exemplar]
By \cite{Obstruction} Proposition~1,
$V_{\mathrm{quad}}$ does not admit a
Gauss--Schwarz triangle interpretation.
Its governing equation has an irregular
singularity at $\infty$, placing it in the
confluent-Heun sector.
It is excluded from all classical special-value
families at \texttt{dps}$=300$--$500$.
The correct uniformization framework, if any
exists, requires a Heun or isomonodromic
theory (\cite{Obstruction}, Conjecture~4).
\end{remark}

\begin{theorem}[Apparent singularity exclusion
  for $V_{\mathrm{quad}}$]
\label{thm:exclusion2}
Let $V_{\mathrm{quad}}$ be the Stokes connection
coefficient $M_{11}$ of the rank-2 ODE
$a(x)\,y''+a'(x)\,y'+c(x)\,y=0$ with
$a(x)=3x^2+x+1$ and $c(x)=-x^2$.
\begin{enumerate}
\item[\textup{(i)}]
  At both finite singularities
  $s_{1,2}=(-1\pm i\sqrt{11})/6$
  (roots of $a$), the indicial equation
  is $\rho^2=0$, giving a double root
  $\rho=0$.
\item[\textup{(ii)}]
  The monodromy around each $s_k$ fixes
  $M_{11}=V_{\mathrm{quad}}$; the effective
  monodromy on the connection coefficient
  is trivial. In particular, $V_{\mathrm{quad}}$
  cannot be identified via CM elliptic periods,
  gamma values, or classical hypergeometric
  special values by any monodromy-based argument.
\end{enumerate}
\end{theorem}

\begin{proof}
The structural key is the identity
$b(x) \equiv a'(x)$, i.e.\
$6x+1 = \frac{d}{dx}(3x^2+x+1)$,
which makes the ODE exact:
$\frac{d}{dx}[a(x)\,y'] + c(x)\,y = 0$.
This forces the indicial exponent
$p_0 = a'(s_k)/a'(s_k) = 1$ and
$q_0 = 0$ at each simple root $s_k$
of $a$, giving indicial equation
$\rho^2 = 0$. The monodromy around
$s_k$ is then unipotent: it fixes
the analytic first solution $y_1$
and acts only on $y_2 = y_1\log\zeta
+ \cdots$, leaving $M_{11}$
invariant.
The full four-step Frobenius computation
and numerical certificate (at
\texttt{dps}$=150$) appear
in~\cite{Vquad}.
\end{proof}

\begin{conjecture}[Master hierarchy
  {\cite{Obstruction}, Conjecture~5}]
\label{conj:master}
Every polynomial continued fraction family in
the present spectral program admits a
well-defined tier within
Definition~\ref{def:tiers}. The tier
determines the correct ambient geometry:
modular ($X(1)$), cuspidal, non-arithmetic
Fuchsian, or irregular Heun.
\end{conjecture}

\begin{theorem}[Tier--stratum correspondence
  {\cite{Obstruction}}]
\label{thm:tiers}
The four obstruction tiers of the JNT paper
correspond to strata as follows:
\begin{center}
\renewcommand{\arraystretch}{1.3}
\begin{tabularx}{\textwidth}{llXl}
\hline
Tier & Geometry & Condition & Stratum \\
\hline
Pre-tier & Wallis collapse &
  $a(1)=0$ or $a(k)=0$ &
  \textbf{Rat} \\
0 & Modular $(\infty,2,3)$ &
  $\lambda=0$, $(\mu,\nu)=(2,3)$, $c=1$ &
  \textbf{Log} (Tier-0) \\
1 & Arithmetic cuspidal &
  Arithmetic origin, Frobenius-log residual &
  \textbf{Log} (Tier-1) \\
2 & Non-arithmetic Fuchsian &
  Gauss $\,_2F_1$ witness, not in
  Takeuchi list &
  \textbf{Des} (Tier-2) \\
3 & Irregular Heun &
  No Gauss witness, irregular
  singularity at $\infty$ &
  \textbf{Des} (Tier-3) \\
\hline
\end{tabularx}
\end{center}
In particular:
\begin{itemize}
\item The \textbf{Log} stratum is the union of
  Tier-0 (modular, generic transcendental limit
  on $X(1)$) and Tier-1 (arithmetic cuspidal,
  logarithmic Frobenius residual).
\item The \textbf{Des} stratum is the union of
  Tier-2 (non-arithmetic Fuchsian, Takeuchi
  obstruction) and Tier-3 (irregular Heun,
  non-Fuchsian).
\item The \textbf{Rat} stratum is pre-tier:
  the Wallis recurrence collapses before any
  hypergeometric structure is reached.
\item The \textbf{Alg} and \textbf{Trans}
  strata are not directly addressed by the
  JNT paper, which operates in the $d=1$
  sector; their tier placement is
  \PLACEHOLDER{open, pending higher-degree
  analysis}.
\end{itemize}
\end{theorem}

% ══════════════════════════════════════════════════════
\section{The Desert Stratum}
\label{sec:desert}
% ══════════════════════════════════════════════════════

\begin{definition}[Desert stratum]
$K\in\mathcal{F}(d,D)$ is Desert if it is
convergent and $\rho(K)=0$: no integer
relation of bounded norm is found against
any basis in
$\mathcal{B}=\{$algebraic, logarithmic,
polylogarithmic, hypergeometric$\}$
at \texttt{dps}$=500$.
\end{definition}

\begin{theorem}[Desert witnesses {\cite{Desert}}]
Degree profiles $(4,3)$ and $(5,3)$ contain
Desert families: PSLQ null at
\texttt{dps}$=500$ against all bases in
$\mathcal{B}$.
\PLACEHOLDER{insert norm bound from accepted
version of 267262713}
\end{theorem}

\begin{remark}[Desert dominance under uniform
  sampling]
\label{rem:desert-dominance}
Under uniform random sampling with $D=4$,
$d=2,\ldots,5$, the Desert stratum accounts
for $957/1000$ families at
\texttt{dps}$=150$.
The Log and Alg strata are not populated by
random sampling at these parameters; known
examples in each stratum were found by
targeted construction.
This suggests
\[
  \mu(\mathrm{Log}\cup\mathrm{Alg})=0
\]
under the uniform measure on
$\mathcal{F}(d,D)$, while both strata are
non-empty and mathematically non-trivial.
\UPDATED{empirical result, 1000-family
classification pass, 2026-04-20}
\end{remark}

\begin{remark}[Genuine ln(2) null result]
A corrected search over 800 families
($d=2,\ldots,5$, $D=4$, seed offset 1500)
using stripped basis $\{1,\ln 2,\ln^2 2\}$
with rationality pre-screen finds zero
genuine $\ln 2$ identities at residual
$<10^{-40}$.
This confirms that $\ln 2$ PCF identities,
if they exist at $D=4$, are not detectable
by random sampling and require targeted
construction.
\UPDATED{confirmed 2026-04-20; see
\cite{Rat} for pre-screening protocol}
\end{remark}

% ══════════════════════════════════════════════════════
\section{Completeness}
\label{sec:completeness}
% ══════════════════════════════════════════════════════

\begin{conjecture}[Completeness]
\label{conj:complete}
Every convergent $K\in\mathcal{F}(d,D)$
falls into exactly one stratum:
\[
  \mathcal{F}(d,D)_{\mathrm{conv}}
  \;=\;
  \mathrm{Rat}\sqcup\mathrm{Log}
  \sqcup\mathrm{Alg}\sqcup\mathrm{Trans}
  \sqcup\mathrm{Des}.
\]
\end{conjecture}

\noindent\textbf{Experimental evidence
(1000-family survey, $d=2,\ldots,6$,
$D=4$):}

\begin{center}
\begin{tabular}{lrrp{5cm}}
\hline
Stratum & Count & \% & Notes \\
\hline
Rat & 43 & 4.3\% &
  Fully characterized \cite{Rat} \\
Log & 0 & 0\% &
  Not found by random sampling;
  examples by construction;
  zero pre-filter proximity to log
  constants in 1500-family directed
  search \cite{Log} \\
Alg & 0 & 0\% &
  Not found by random sampling;
  examples by construction \cite{Alg,Spectral} \\
Trans & 0 & 0\% &
  No identified hits at \texttt{dps}$=150$ \\
Des & 957 & 95.7\% &
  Dominant stratum \cite{Desert,Khinchin} \\
\textbf{Both (artifact)} & \textbf{41$\to$0} & — &
  All confirmed rational by pre-screen;
  see Remark~\ref{rem:rat-contamination} \\
\textbf{Unclassified} & \textbf{0} & \textbf{0\%} &
  Completeness holds empirically \\
\hline
\end{tabular}
\end{center}

\begin{remark}[Both-stratum artifact]
\label{rem:rat-contamination}
The initial classification flagged 41 families
as hitting multiple bases (``Both''). Rationality
pre-screening via $[1,K]$ PSLQ confirmed all
41 have $K\in\mathbb{Q}$; the multi-basis hits
were spurious expressions of rational limits.
The 11 distinct rational values include
$K\in\{-2,-1,0,1,2,3,4\}$ and two simple
fractions ($29/8$, $21/4$), with $K=1$ most
frequent ($9$ families across $d=2,\ldots,6$).
\UPDATED{diagnosed 2026-04-20}
\end{remark}

\begin{remark}[Relation to F1]
Conjecture~\ref{conj:complete} is equivalent
to the \emph{generalized degree-$d$ dichotomy}
(F1): for any degree $d\geq 2$, every
convergent PCF has an arithmetic type
detectable by $\mathrm{Spec}(K)$ at finite
precision.
\end{remark}

% ══════════════════════════════════════════════════════
\section{Formal Verification Status}
\label{sec:lean}
% ══════════════════════════════════════════════════════

The Wallis--PCF family has been partially
formalized in Lean~4 (repository:
\href{https://github.com/papanokechi/siarc-lean4}{%
github.com/papanokechi/siarc-lean4}).

\begin{itemize}
\item \textbf{36 proof targets} discharged
  in \texttt{proof\_status.json}
  (iterations 12--13).
\item \textbf{Poincar\'e characteristic lemma}
  transcribed from \cite{Obstruction}.
\item \textbf{CasoratianPi4.lean} and
  \textbf{ShiftConsistency.lean} repaired
  (iteration 13).
\item \textbf{WallisFamily.lean}: 
  \DONE{lake build clean, 8267 jobs,
  0 errors, 0 sorry, 0 warnings.
  Fix: \texttt{open scoped Topology} +
  \texttt{push\_cast} cast normalization.
  Completed iteration 15, 2026-04-20.}
\item \textbf{Lemma K} (Kloosterman sum bound
  at conductor $N=24$): still blocked;
  key theoretical blocker for G-01
  universal law conjecture $k\geq 5$.
\end{itemize}

% ══════════════════════════════════════════════════════
\section{Open Problems}
\label{sec:open}
% ══════════════════════════════════════════════════════

\begin{enumerate}
\item \textbf{(F1)} Prove
  Conjecture~\ref{conj:complete} for $d=2$,
  $D=4$.

\item \textbf{(F2)} Identify
  $V_{\mathrm{quad}}\approx 1.19737\ldots$
  The governing ODE has the exact form
  $a(x)y''+a'(x)y'+c(x)y=0$ with
  $a(x)=3x^2+x+1$, $c(x)=-x^2$; both
  finite singularities are apparent
  (trivial monodromy, indicial exponents
  $\{0,0\}$) by the algebraic identity
  $b(x)\equiv a'(x)$. The CM discriminant
  $-11$ is $\mathrm{disc}(a)=1-12=-11$,
  not a deep CM structure. The constant
  equals the Stokes connection coefficient
  $M_{11}$ of the irregular singularity;
  its isomonodromic deformation is
  PIII($D_6$) with $\alpha=1/6$,
  $\beta=\gamma=0$, $\delta=-1/2$.
  The Stokes constant $S\approx 0.43770528$
  and connection parameter
  $\sigma_{\mathrm{conn}}\approx 0.060877$
  are computed but unidentified in closed
  form; identification requires the
  PIII($D_6$) $\tau$-function at the CM
  accessory parameter $q=(5+i\sqrt{11})/54$.
  See \cite{Vquad} for the complete analysis.
  \UPDATED{submitted to Nonlinearity
  NON-110708, 2026-04-21}

\item \textbf{(F3)} Give a computable
  definition of the Arithmetic Valuation
  Operator (AVO) extending $\tau$.

\item \textbf{(F4)} Determine whether the
  Desert stratum contains provably
  transcendental elements.

\item \textbf{(F5)} Characterize the Alg
  stratum for $d\geq 3$: does every
  degree-$d$ algebraic PCF limit arise as a
  period of an arithmetic Fuchsian group?

\item \textbf{(F6, new)} Explain the Desert
  dominance: prove that
  $\mu(\mathrm{Log}\cup\mathrm{Alg})=0$
  under the uniform measure on
  $\mathcal{F}(d,D)$, or exhibit a
  constructive density lower bound.

\item \textbf{(F7, new)} Determine the
  asymptotic Desert rate as $D\to\infty$
  and $d\to\infty$ separately.
\end{enumerate}

% ══════════════════════════════════════════════════════
\section*{Acknowledgments}
% ══════════════════════════════════════════════════════

Results were discovered through the SIARC
multi-model relay pipeline across iterations
1--\PLACEHOLDER{final iteration number}.
All computations are reproducible via the
supplementary repositories.

\smallskip
\noindent\textit{AI assistance disclosure.}
The preparation of this manuscript involved
assistance from Claude
(\texttt{claude-sonnet-4-6}, Anthropic) for
\LaTeX{} editing, proof review, and
computational interpretation, and from
GitHub Copilot (Microsoft) for autonomous
script execution within the SIARC relay chain.
All mathematical results, proofs, and
numerical verifications were conceived,
directed, and validated by the author.
The AI was used as a writing and research
tool, not as a source of original mathematics.

% ══════════════════════════════════════════════════════
\begin{thebibliography}{99}

\bibitem{Rat}
  papanokechi,
  \textit{Trivial Rational Contamination in
  PSLQ-Based Polynomial Continued Fraction
  Searches: Diagnosis, Correction, and a
  Pre-Screening Protocol},
  submitted to Experimental Mathematics,
  submission 261945835 (2026-04-20).

\bibitem{Log}
  papanokechi,
  \textit{Polynomial Continued Fractions:
  a Proved Logarithmic Ladder, a $4/\pi$
  Casoratian Identity, and 482 Irrational
  Constants},
  \PLACEHOLDER{Experimental Mathematics,
  submission 266999523, under review;
  Zenodo DOI: 10.5281/zenodo.19491767}.

\bibitem{Alg}
  papanokechi,
  \textit{Arithmetic Sector Bifurcation in
  Quadratic Continued Fractions: A Formalized
  Study of the Wallis--PCF Family},
  \PLACEHOLDER{Experimental Mathematics,
  submission 266535114, under review}.

\bibitem{Obstruction}
  papanokechi,
  \textit{PSL$_2(\mathbb{Z})$ and a Four-Tier
  Obstruction Hierarchy in the Spectral Theory
  of Polynomial Continued Fractions},
  \PLACEHOLDER{Journal of Number Theory,
  submission JNTH-D-26-00480, under review}.

\bibitem{Spectral}
  papanokechi,
  \textit{Spectral Classes of Polynomial
  Continued Fractions: A Hypergeometric
  Framework and Arithmetic Obstructions},
  \PLACEHOLDER{Experimental Mathematics,
  submission 268704746, under review}.

\bibitem{Desert}
  papanokechi,
  \textit{A Computational Verification of the
  Polynomial Continued Fraction Desert at
  Degree Profiles $(4,3)$ and $(5,3)$},
  \PLACEHOLDER{Experimental Mathematics,
  submission 267262713, under review}.

\bibitem{Khinchin}
  papanokechi,
  \textit{A Precision-Controlled Null Result
  for a Khinchin-Signature PSLQ Family},
  \PLACEHOLDER{Experimental Mathematics,
  submission 264514392, under review}.

\bibitem{Vquad}
  papanokechi,
  \textit{Painlev\'e~III($D_6$) and resurgence
  for a constant from a quadratic polynomial
  continued fraction},
  submitted to Nonlinearity,
  submission NON-110708 (2026-04-21).

\bibitem{Takeuchi}
  K.~Takeuchi,
  \textit{Arithmetic triangle groups},
  J.\ Math.\ Soc.\ Japan \textbf{29} (1977),
  no.~1, 91--106.

\bibitem{Yoshida}
  M.~Yoshida,
  \textit{Fuchsian Differential Equations:
  With Special Emphasis on the Gauss--Schwarz
  Theory},
  Vieweg, Braunschweig, 1987.

\bibitem{ClarkVoight}
  P.~L.~Clark and J.~Voight,
  \textit{Algebraic curves uniformized by
  congruence subgroups of triangle groups},
  Trans.\ Amer.\ Math.\ Soc.\
  \textbf{363} (2011), no.~1, 33--60.

\bibitem{RamanujanMachine}
  Raayoni, G.\ et al.,
  \textit{Generating conjectures on fundamental
  constants with the Ramanujan Machine},
  Nature \textbf{590} (2021), 67--73.

\bibitem{LorentzenWaadeland}
  Lorentzen, L.\ and Waadeland, H.,
  \textit{Continued Fractions with
  Applications}, North-Holland, 1992.

\bibitem{aeal-paper7}
  papanokechi,
  \textit{AI Behavior Science --- 
  Founding Territory Paper},
  Zenodo,
  \href{https://doi.org/10.5281/zenodo.19562751}{%
  DOI: 10.5281/zenodo.19562751} (2026).

\bibitem{aeal-paper8}
  papanokechi,
  \textit{AEAL: Agent Epistemic
  Accountability Layer},
  Zenodo,
  \href{https://doi.org/10.5281/zenodo.19565086}{%
  DOI: 10.5281/zenodo.19565086} (2026).

\end{thebibliography}

\end{document}